We have presented \texttt{LimberCloud}, a tensorised reformulation of the Limber projection in which every angular power spectrum of the \(3\times2\,\mathrm{pt}\) data vector is expressed as a bilinear contraction, \(C_{uv}^{ab}(\ell) = \sum_{ij} B_{uv}^{ij}(\ell)\,\phi^a_i\,\phi^b_j\) (Equation~\ref{eq:5.6}). The basis tensor \(B_{uv}^{ij}(\ell)\) encapsulates all dependence on cosmological and astrophysical parameters, while the discretised tomographic redshift distributions \(\phi^a_i\) encode the survey-specific information. This factorisation is obtained by piecewise-linearising the radial kernels and the three-dimensional power spectrum on a discrete redshift grid, reducing each interval contribution to one of 21 structurally distinct closed-form integrals that populate the basis tensor without numerical quadrature.

Validated against \texttt{CCL} for the full LSST Y1 \(3\times2\,\mathrm{pt}\) configuration---encompassing cosmic shear, galaxy--galaxy lensing, and galaxy clustering across all tomographic bin combinations---the framework achieves sub-per-cent accuracy over the multipole range \(20 \le \ell \le 2000\), with typical residuals one to two orders of magnitude below the expected survey variance.

Benchmark timings demonstrate that the \texttt{JAX--GPU} implementation reduces the cost of the projection step to approximately \(3\,\mathrm{ms}\) per evaluation for the full \(3\times2\,\mathrm{pt}\) data vector, more than three orders of magnitude faster than the conventional \texttt{CCL} pipeline. The cosmology stage---computing \(\chi(z)\) and \(P_{\delta\delta}(k,z)\) via \texttt{CAMB}---dominates the total wall-clock budget but is amenable to replacement by existing power-spectrum emulators. The tensorised structure further enables end-to-end emulation of the basis tensor itself: training a lightweight surrogate on \(\sim\!10^6\) parameter vectors is feasible within \(\mathcal{O}(1)\)~day on a modest GPU cluster and would reduce the per-evaluation cost to a neural-network forward pass plus the inexpensive projection contraction, yielding an expected speed-up of three to four orders of magnitude over the conventional approach.

The bilinear dependence of the angular power spectra on the redshift distributions yields an exact polynomial expansion of the data vector in the nuisance parameters (Equation~\ref{eq:6.3}), terminating at second order with no truncation error. Under a Gaussian prior, the linear-order contribution permits analytic marginalisation by completing the square, inflating the data covariance by the cosmology-dependent term \(\boldsymbol{\Gamma} \, \boldsymbol{\Pi} \, \boldsymbol{\Gamma}^{\mathrm{T}}\) (Equation~\ref{eq:6.8}). Higher-order corrections are accessible from first principles via Isserlis' theorem and are expected to be small when photometric calibration priors are well constrained. This treatment avoids the finite-difference derivatives and fixed-cosmology approximations employed in existing analytic marginalisation schemes, providing a controlled framework for quantifying the adequacy of the linear-order approximation.

The framework extends naturally to non-flat geometries, to additional projected observables---including CMB lensing convergence, the integrated Sachs--Wolfe effect, and redshift-space distortions---and to beyond-Limber corrections via spherical-Bessel or FFTLog decompositions, all of which preserve the bilinear factorisation and its associated emulation and marginalisation infrastructure. Follow-up work will present end-to-end emulation within a differentiable inference pipeline and numerical validation of the analytic marginalisation formalism.
